
 To solve this exercise, we remember that if a variable $x$ can take any real value, it can be replaced by two non negative artificial variables denoted by $x^+$ and $x^-$, such that $x = x^+-x^-$. We also recall that the absolute value of $x$ is defined as
$$ |x| = \left\{ 
\begin{array}{rl}
 x & \mbox{if } x \geq 0, \\ 
 -x & \mbox{if } x < 0.
\end{array}
\right.$$

In our case, we have that
$$ |x_1 - x_2| = \left\{ 
\begin{array}{rl}
 x_1 - x_2 & \mbox{if } x_1 \geq x_2, \\ 
x_2 - x_1 & \mbox{if } x_1 <x_2.
\end{array}
\right.$$
Since $x_1$ and $x_2$ are non negative real numbers, the difference $x_1 - x_2$ can take any real value.

Let us study the absolute value by examining the two cases:
\begin{itemize}
\item If $x_1 - x_2 \geq 0$, we can define the non negative quantity
\[
y^+ = \left\{
\begin{aligned}
  x_1 - x_2, & \text{ if } x_1-x_2 > 0, \\
  0,     & \text{ otherwise}.
\end{aligned}
\right.
\]

\item If $x_1 - x_2 < 0$, we can define the positive quantity
\[
y^- = \left\{
\begin{aligned}
  x_2 - x_1, & \text{ if } x_1-x_2 < 0, \\
  0,     & \text{ otherwise}.
\end{aligned}
\right.
\]
\end{itemize}

Consequently, the absolute value of the difference can then be written as
$$ | x_1 - x_2 | = y^+ + y^-. $$

Furthermore, if we impose that $y^+ \geq 0$ and $y^- \geq 0$, we have that our minimization problem can be written as
$$\min_{x_1,x_2,y^+,y^-} (y^+ + y^-) $$
subject to
\begin{align*}
y^+ & \geq x_1-x_2,\\
y^- & \geq x_2-x_1, \\
x_1 & \geq 0,\\
x_2 & \geq 0,\\
y^+ & \geq 0,\\
y^- & \geq 0.
\end{align*}

Note that this formulation is not strictly equivalent to the original one for
all feasible solutions. For instance, the feasible solution $x_1=0$,
$x_2=0$, $y^+=1$, $y^-=1$ has objective value 2 in the transformed
problem, and 0 in the original one.

But it is equivalent at the optimal solution. Denote the optimal
solution of the original problem $(x_1^*,x_2^*)$.
\begin{itemize}
\item If $x_1^*-x_2^* > 0$, the lowest possible value for $y^+$,
  denoted by $(y^+)^*$, is
  $x_1^*-x_2^*$ (because of the constraint $y^+  \geq x_1-x_2$), and
  the lowest possible value for $y^-$, denoted by $(y^-)^*$, is $0$
  (because of the constraint $y^- \geq 0$). Therefore,
  the objective function of the transformed problem at the optimal
  solution is
  \[
(y^+)^* + (y^-)^* = x_1^*-x_2^* + 0 = |x_1^*-x_2^*|.
  \]
\item If $x_1^*-x_2^* < 0$, for similar reasons, we have
  $(y^+)^*=0$ and $(y^-)^*=x_2^*-x_1^*$. Therefore,
  the objective function of the transformed problem at the optimal
  solution is
  \[
(y^+)^* + (y^-)^* = 0 +x_2^*-x_1^* = |x_1^*-x_2^*|.
\]
\item If $x_1^*-x_2^* = 0$, we have
  $(y^+)^*=0$ and $(y^-)^*=0$. Therefore,
  the objective function of the transformed problem at the optimal
  solution is
  \[
(y^+)^* + (y^-)^* = 0 + 0 = |x_1^*-x_2^*|.
\]
\end{itemize}




